% !TeX root = proyecto.tex

% CAPÍTULO 8: RESULTADOS Y EVALUACIÓN

\chapter{Resultados y evaluación}
\label{cap:resultados}

Este capítulo expone de forma analítica y cuantitativa el desempeño técnico del sistema
integrado al cierre del trabajo terminal II. Se demuestra la superación de los objetivos
cuantitativos establecidos al inicio del TT2 mediante los resultados de una corrida de
producción completa sobre el corpus de validación final.

\section{Resultados de la recolección y deduplicación}
\label{sec:res_recoleccion}

\subsection{Volumen de ingesta}
\label{subsec:res_volumen}

La migración del esquema de recolección del TT1 (10 feeds RSS con 281 noticias acumuladas
en múltiples semanas) hacia la arquitectura híbrida del TT2 (45 feeds RSS nacionales y de
género, consultas a Google News y cobertura complementaria mediante \texttt{StealthSession})
produjo un incremento sustancial y medido en la capacidad de ingesta. En el ciclo de validación
final, el módulo recolector procesó al rededor de \textbf{900 noticias brutas} en un único bloque de
ejecución continuo, de las cuales la cascada de deduplicación generó \textbf{634 noticias únicas}
para análisis semántico.

\subsection{Cascada de deduplicación}
\label{subsec:res_dedup}

La deduplicación en cascada de cuatro fases eliminó \textbf{213 duplicados} (25\% del
corpus bruto) distribuidos por etapa de la siguiente forma:

\begin{table}[H]
\centering
\caption{Distribución de duplicados eliminados por fase de la cascada de deduplicación}
\label{tab:res_dedup}
\renewcommand{\arraystretch}{1.3}
\begin{tabular}{|c|p{4.5cm}|c|c|}
\hline
\rowcolor{black}
\textcolor{white}{\textbf{Fase}} &
\textcolor{white}{\textbf{Algoritmo}} &
\textcolor{white}{\textbf{Duplicados eliminados}} &
\textcolor{white}{\textbf{\% acumulado}} \\
\hline
1 & Hash MD5 exacto & 89 & 10.5\% \\
\hline
2 & Similitud de Jaccard ($J \geq 0.70$) & 61 & 17.7\% \\
\hline
3 & SimHash (distancia Hamming $\leq 6$ bits) & 43 & 22.8\% \\
\hline
4 & Similitud coseno TF-IDF ($\cos \geq 0.85$) & 20 & 25.1\% \\
\hline
\textbf{Total} & & \textbf{213} & \textbf{25\%} \\
\hline
\end{tabular}
\end{table}

Este resultado confirma la hipótesis de diseño pues la mayor parte de la redundancia periodística
(un 10.5\%) corresponde a re-publicaciones literales identificables mediante hash MD5, mientras
que un 14.6\% adicional corresponde a reescrituras con vocabulario diferente que solo pueden
detectarse mediante métricas graduales de similitud textual (Jaccard, SimHash, coseno TF-IDF).

\section{Resultados de la clasificación semantica escalonada} 
\label{sec:res_deteccion}

\subsection{Distribución del embudo analítico}
\label{subsec:res_embudo}

El corpus de 634 noticias únicas fue procesado secuencialmente por las cuatro capas del
sistema de clasificación escalonada. La siguiente tabla muestra el volumen
de noticias en cada etapa del embudo.

\begin{table}[H]
\centering
\caption{Distribución de noticia por etapa en la clasificación escalonada}
\label{tab:res_embudo}
\renewcommand{\arraystretch}{1.3}
\begin{tabular}{|c|p{5.5cm}|c|c|}
\hline
\rowcolor{black}
\textcolor{white}{\textbf{Capa}} &
\textcolor{white}{\textbf{Etapa}} &
\textcolor{white}{\textbf{Noticias}} &
\textcolor{white}{\textbf{\% del corpus}} \\
\hline
Entrada & Corpus deduplicado & 634 & 100.0\% \\
\hline
Capa 0 & Escudo léxico (dominios ajenos eliminados) & 409 & 64.5\% \\
\hline
Capa 1 & Escudo suave (dominios NNA/violencia, con penalización) & 225 & 35.5\% \\
\hline
Capa 2 & Inferencia BETO (score $> 0.50$) & 128 & 20.2\% \\
\hline
Capa 3 & Post-filtro de validación aplicado & 128 & 20.2\% \\
\hline
\textbf{Alta} & \textbf{Casos detectados de orfandad} & \textbf{27} & \textbf{4.4\%} \\
\hline
\textbf{Media} & \textbf{Contextos de vulnerabilidad colateral} & \textbf{101} & \textbf{15.9\%} \\
\hline
\end{tabular}
\end{table}


\subsection{Comportamiento por modo de inferencia}
\label{subsec:res_inferencia}

El sistema opero en modo \textit{Zero-Shot} para el corpus de validación inicial durante el prototipo 3, y en modo \textit{Finetuned}
una vez disponible el clasificador ajustado en el prototipo 4. La diferencia operativa entre ambos modos fue muy notable.

\begin{itemize}
    \item \textbf{Modo \textit{Zero-Shot} (P3):} Los scores BETO se concentraban en el
          rango estrecho $[0.40, 0.70]$ lo que generaba el efecto péndulo documentado en el
          capítulo de desarrollo. Las reglas binarias del post-filtro degradaban sistemáticamente
          casos genuinos cuya redacción periodística no coincidía con los patrones esperados.

    \item \textbf{Modo \textit{Finetuned} (P4):} El clasificador ajustado produjo scores
          dispersos en todo el rango $[0, 1]$ lo que permitia que el escudo suave y el
          post-filtro operaran sobre una distribución de probabilidad rica y discriminativa.
          El efecto péndulo fue eliminado, y las 27 noticias de Alta clasificación
          correspondian a casos de orfandad verificables manualmente.
\end{itemize}

\section{Resultados del Agrupamiento Semántico (BERTopic)}
\label{sec:res_agrupamiento}


El flujo de BERTopic (UMAP + HDBSCAN + c-TF-IDF) procesó las 634 noticias deduplicadas
y generó \textbf{15 tópicos semánticos} con alta relacion semantica entre si. El porcentaje inicial
de outliers (no asociados a ningún tópico) fue del \textbf{81.4\%} el cual fue reducido a un \textbf{55\%}
mediante la estrategia de reducción multi-etapa de tres fases (probabilidades blandas de HDBSCAN,
distribuciones c-TF-IDF y cercanía de embeddings con umbral 0.3).

La calidad del agrupamiento del TT2 representa una mejora cualitativa respecto al TT1.
Mientras K-Means (P1) producia un \textit{Silhouette Score} de \textbf{0.23} operando
sobre vectores TF-IDF léxicos, los tópicos BERTopic del TT2 operan sobre embeddings BETO
de 768 dimensiones reducidos a 5D por UMAP y lo que permitia agrupar noticias sobre el mismo
caso criminal publicadas por medios con vocabularios distintos. Entre los tópicos identificados se
distinguen categorías temáticas como casos de orfandad con intervención institucional
del DIF, disputas legales por la custodia de menores, cobertura de feminicidios con
NNA testigos y manifestaciones de colectivos de mujeres.

\section{Resultados del Módulo de Investigación Profunda (\textit{DeepInvestigator})}
\label{sec:res_investigador}

El módulo \texttt{DeepInvestigator} fue activado sobre una muestra de
\textbf{15 noticias de Alta clasificación} durante las pruebas de validación del
prototipo 4. En todos los casos el módulo completó satisfactoriamente los cuatro
pasos del flujo de investigación (generación de query con Gemini Flash, búsqueda
en DuckDuckGo, scraping complementario mediante \texttt{StealthSession} y síntesis
estructurada en JSON) con un tiempo promedio de \textbf{47 segundos por caso}.

La ficha de caso generada para cada noticia incluye los siete campos estructurados
(\texttt{ubicacion}, \texttt{victimas}, \texttt{ninos\_afectados}, \texttt{edades},
\texttt{situacion\_actual}, \texttt{medios\_contacto} y \texttt{resumen}),
con el campo \texttt{medios\_contacto} siendo el de mayor valor operativo pues
Gemini retorna el DIF municipal y la Fiscalía competente según la ubicación
del caso aun incluso cuando el texto de la noticia no los menciona explícitamente.

\section{Evaluación comparativa TT1/TT2 y cumplimiento de metas}
\label{sec:res_comparativa}

La siguiente tabla consolida el cumplimiento de los objetivos cuantitativos
definidos al inicio del TT2 frente a las métricas del TT1 (prototipo 2) y los
resultados finales del prototipo 4.

\begin{table}[H]
\centering
\caption{Evaluación comparativa de métricas TT1/P2 vs. metas TT2 vs. resultados P4}
\label{tab:res_comparativa}
\renewcommand{\arraystretch}{1.4}
\begin{tabular}{|p{4.0cm}|c|c|c|c|}
\hline
\rowcolor{black}
\textcolor{white}{\textbf{Métrica}} &
\textcolor{white}{\textbf{TT1 (P2)}} &
\textcolor{white}{\textbf{Meta TT2}} &
\textcolor{white}{\textbf{Resultado P4}} &
\textcolor{white}{\textbf{Cumplimiento}} \\
\hline
\textbf{Falsos positivos} & 60\% & $< 30\%$ & $\sim 20\%$ (estimado) & \checkmark \\
\hline
\textbf{Fuentes monitoreadas} & 10 RSS & 45 RSS + Google News & 45 RSS + Google News + Stealth & \checkmark \\
\hline
\textbf{Cobertura temporal} & 6 meses & Continua & Operacional & \checkmark \\
\hline
\textbf{Persistencia} & CSV plano & PostgreSQL 16 & PostgreSQL 16 (ACID) & \checkmark \\
\hline
\textbf{Motor de clasificación} & RegEx (72 patrones) & BETO fine-tuning & BETO Finetuned + Escalonado & \checkmark \\
\hline
\textbf{Volumen de ingesta} & 281 noticias & $> 500$ únicas & 634 únicas & \checkmark \\
\hline
\end{tabular}
\end{table}

\begin{figure}[htbp]
    \centering
    % \includegraphics[width=0.90\textwidth]{images/lista_noticias_alta}
    \rule{0.90\textwidth}{5.5cm}
    \caption{Visualización jerárquica de noticias de Alta relevancia en el Dashboard, con etiquetas
    de clasificación y porcentaje de score BETO por registro.}
    \label{fig:lista_noticias_alta}
\end{figure}

La interfaz del Dashboard con etiquetas visuales (clasificación Alta/Media/Baja
y score semántico) permiten al analista realizar un barrido visual inmediato sobre
el corpus clasificado. Lo que en el TT1 representaba decenas de horas para una revisión
manual fue optimizado en el TT2 a una revisión semi-automatizada que exige apenas
unos pocos segundos por caso, antes de activar el módulo de investigación profunda para los
casos que el analista determina que merecen una intervención más exhaustiva.

